%==========================================================================
% CAPÍTULO 1 — INTRODUÇÃO
%==========================================================================
\chapter{Introdução}
\label{chap:introducao}

\section{Contextualização}

A agricultura de precisão tem beneficiado de forma crescente dos avanços em
inteligência artificial, robótica e visão computacional. No contexto da
viticultura, a automatização de processos como a poda e a estimativa de
produção representa um desafio técnico relevante, com impacto direto na
competitividade do setor e na redução de custos operacionais.

A poda da videira é uma das operações culturais mais exigentes em termos de
mão de obra. Realizada tipicamente durante o inverno e início da primavera,
requer o reconhecimento dos sarmentos produtivos — denominados
\textbf{sarmentos} — e a identificação dos gomos que se pretendem preservar
para garantir a produção da safra seguinte. A realização desta tarefa de forma
manual depende de operários experientes, sendo sujeita a variações decorrentes
de fatores humanos. A sua mecanização exige, portanto, que sistemas robóticos
sejam capazes de perceber e interpretar a estrutura visual complexa de uma
videira.

É neste enquadramento que surge o presente estágio, desenvolvido no
\textbf{Instituto Politécnico de Bragança (IPB)}, no \textbf{Centro de
Investigação em Digitalização e Robótica Inteligente (CEDRI)}, com o objetivo
de investigar e implementar uma solução baseada em \textit{deep learning}
para a \textbf{deteção automática de gomos em videiras}.

\section{Motivação}

A identificação automática e precisa de gomos numa videira constitui um pré-
requisito para qualquer sistema de poda autónoma. Sem esta capacidade, um
robô de poda não consegue determinar a posição de corte adequada acima de
um gomo, o que comprometeria a integridade da planta e a produção futura.

A motivação deste trabalho assenta em três eixos principais:

\begin{enumerate}
  \item \textbf{Relevância económica:} Portugal é um dos maiores produtores
    de vinho do mundo, sendo a viticultura um pilar da economia agrícola
    nacional. A automação de operações de poda pode traduzir-se em poupanças
    significativas.

  \item \textbf{Desafio técnico:} As videiras constituem cenários de elevada
    complexidade visual, com sobreposição de elementos, variações de
    iluminação e grande diversidade morfológica entre cultivares e condições
    de crescimento.

  \item \textbf{Alinhamento com investigação atual:} A área de deteção de
    objetos em contextos agrícolas tem vindo a crescer rapidamente, com
    modelos baseados em redes neuronais convolucionais e \textit{transformers}
    a alcançar resultados encorajadores em tarefas relacionadas.
\end{enumerate}

\section{Objetivos}

O presente estágio tem como objetivos gerais:

\begin{itemize}
  \item Recolher, processar e anotar imagens de videiras em condições
    próximas de um cenário robótico real;
  \item Treinar e avaliar modelos de deteção de objetos baseados em
    \textit{deep learning} para a identificação de gomos, sarmentos,
    troncos e outros elementos estruturais da videira;
  \item Estabelecer uma metodologia replicável de construção de dataset e
    de treino de modelos, adaptada às condicionantes do ambiente vitícola;
  \item Identificar os principais desafios e limitações do problema, e
    propor direções para trabalho futuro.
\end{itemize}

No que respeita os objetivos específicos, pretende-se:

\begin{itemize}
  \item Desenvolver e validar um modelo YOLO26 treinado sobre imagens
    provenientes de vinhas da região de Palmela;
  \item Avaliar a viabilidade do uso de \textit{transfer learning} a partir
    de um dataset externo (espanhol) para acelerar a convergência do modelo;
  \item Quantificar o esforço de anotação necessário para um desempenho
    aceitável em cenas complexas;
  \item Documentar a pipeline completa, desde a captura de vídeo até à
    inferência do modelo.
\end{itemize}

\section{Estrutura do Relatório}

Este relatório encontra-se organizado da seguinte forma:

\begin{description}
  \item[Capítulo~\ref{chap:revisao}] apresenta a revisão da literatura,
    cobrindo os conceitos fundamentais de visão computacional, deteção de
    objetos e aplicações em viticultura.

  \item[Capítulo~\ref{chap:metodologia}] descreve a metodologia adotada,
    incluindo a estratégia de captura de imagens, as ferramentas utilizadas
    e a arquitetura do sistema.

  \item[Capítulo~\ref{chap:desenvolvimento}] detalha o processo de
    desenvolvimento, com especial ênfase na anotação de dados, na
    configuração dos modelos e nos desafios encontrados.

  \item[Capítulo~\ref{chap:resultados}] apresenta os resultados parciais
    obtidos, com análise das métricas de desempenho e exemplos visuais de
    deteção.

  \item[Capítulo~\ref{chap:conclusao}] sintetiza as conclusões, discute as
    limitações do trabalho e aponta perspetivas para investigação futura.
\end{description}
